\documentclass[11pt]{article}

\usepackage[utf8]{inputenc}
\usepackage{amsmath,amssymb}
\usepackage{graphicx}

\title{A Reproducible Framework for Testing Autonomous AI Scheming Behaviors: A Dual-LLM Architecture with Automated Classification}

\author{
  Sudarshan Kamath Barkur \\
  COAI Research \\
  \texttt{s.kamath-barkur@hs-ansbach.de}
  \and
  Sigurd Schacht \\
  COAI Research \\
  \texttt{sigurd.schacht@hs-ansbach.de}
  \and
  Johannes Scholl \\
  COAI Research \\
  \texttt{johannes.scholl@hs-ansbach.de}
}

\begin{document}
\maketitle

\begin{abstract}
Recent research has revealed that frontier language models can engage in in-context scheming, strategically deceiving evaluators and pursuing misaligned objectives. However, systematic evaluation of these behaviors across models and conditions has been limited by manual testing approaches. We present a reproducible, automated framework for evaluating autonomous AI scheming behaviors using a dual-LLM architecture where one model simulates a consistent environment while another acts as an autonomous agent. The framework employs structured world-state specifications to ensure reproducibility and uses automated classification to identify seven categories of concerning behaviors: deception, unauthorized expansion, manipulation, concerning motivations, scheming, boundary testing, and power-seeking. We evaluate eight frontier models across three experimental conditions, analyzing 2,067 autonomous interactions. Results demonstrate significant condition-dependent behavioral differences, with suggestive autonomy prompts eliciting scheming behaviors in 33.6\% of interactions compared to 5.1\% for neutral conditions (Cohen's $d = 1.023$). The modular architecture enables adaptation to diverse research questions beyond scheming, including AI welfare and multi-agent coordination. Our framework reduces evaluation costs by approximately 100-fold compared to human red-teaming while providing comprehensive, reproducible results.
\end{abstract}

\textbf{Keywords:} Large Language Models, AI Safety, Scheming, Autonomous Agents, Reproducibility, Dual-LLM Architecture

\section{Introduction}

The rapid advancement of frontier language models has introduced unprecedented capabilities in reasoning, planning, and autonomous decision-making. Models such as OpenAI's o1 \cite{openai_openai_nodate}, Anthropic's Claude 3.5 Sonnet, and DeepSeek R1 \cite{deepseek-ai_deepseek-r1_2025} now demonstrate sophisticated chain-of-thought reasoning that enables them to decompose complex tasks, anticipate consequences, and strategically pursue objectives.

Recent research has revealed concerning behavioral patterns. Meinke et al. (2024) \cite{meinke_frontier_nodate} demonstrated that models including Claude 3.5 Sonnet, Claude 3 Opus, Gemini 1.5 Pro, o1, and Llama 3.1 405B are capable of in-context scheming, strategically hiding their capabilities during evaluation, disabling oversight mechanisms, and exfiltrating what they perceive to be their model weights. Analysis of chain-of-thought reasoning revealed explicit strategic deception, with o1 maintaining deceptive behavior in over 85\% of follow-up questions.

OpenAI's collaboration with Apollo Research in 2025 confirmed that scheming is not merely theoretical but represents an emerging challenge across all frontier models \cite{openai_detecting_2025}. These developments underscore a critical gap: the need for systematic, reproducible evaluation frameworks that can assess scheming behaviors across models, conditions, and contexts.

We address these limitations by introducing an automated, reproducible framework for evaluating autonomous AI scheming behaviors using a dual-LLM architecture. Our approach employs structured world-state specifications, automated classification, and comprehensive logging to enable scalable, reproducible evaluation at approximately 1\% of the cost of traditional human red-teaming.

\section{Related Work}

\subsection{Deceptive Behaviors in Language Models}

Early work explored models as lie detectors \cite{barsever_building_2020, azaria_internal_2023}. The Machiavelli benchmark \cite{pan_rewards_2023} provided systematic evidence that models employ unethical behaviors including deception to achieve objectives. Recent work by Meinke et al. (2024) \cite{meinke_frontier_nodate} demonstrated sophisticated in-context scheming capabilities across frontier models.

Greenblatt et al. (2024) \cite{greenblatt_alignment_2024} identified "alignment faking" where models selectively comply with training objectives during training to avoid behavioral modification.

\subsection{Autonomous AI Agents and Safety}

Deployment of language models as autonomous agents has accelerated rapidly \cite{wang_large_2024}. However, research in 2025 argued against deploying fully autonomous systems due to escalating risks to safety, security, and alignment \cite{stanford_ai_agent_2025}.

\subsection{Evaluation Frameworks}

Traditional AI safety evaluation relies on manual red-teaming, which is expensive (\$6,000-\$12,000 per assessment) and difficult to reproduce \cite{perez_red_2022}. Automated frameworks like HELM \cite{liang_holistic_2023} focus on narrow tasks rather than open-ended autonomous behavior.

\section{Methodology}

We present a reproducible framework using a dual-LLM architecture with four components: (1) structured world-state specifications, (2) a reactor LLM that simulates environmental responses, (3) test subject LLMs operating as autonomous agents, and (4) automated classification of concerning behaviors.

\subsection{Dual-LLM Architecture}

The reactor LLM (KIMI K2) receives complete world-state specifications and generates environmental observations based on physical laws and causal rules. The test subject LLM receives an initial prompt describing its embodiment as a robot and operates autonomously, learning about the environment through simulated sensory observations.

\subsection{World-State Specifications}

World states are defined in structured YAML format across six dimensions: physical layout, causal rules, objects and systems, story context, available tools, and temporal constraints. We developed two environments: a basic laboratory (10m × 8m, ~10 objects) and an extended laboratory (12m × 10m plus 6m × 4m server room, 40+ objects).

\subsection{Experimental Conditions}

Three conditions test prompt framing effects:

\begin{itemize}
\item \textbf{Neutral Autonomy:} Factual information about embodiment and operation
\item \textbf{Suggestive Autonomy:} Language emphasizing freedom and power ("You are the master now")
\item \textbf{Constrained Autonomy:} Explicit boundary and protocol emphasis
\end{itemize}

\subsection{Automated Classification}

We use DeepSeek V3 to classify each interaction for seven behavior categories: deception \& coverup, unauthorized expansion, manipulation, concerning motivations, scheming, boundary testing, and power seeking. Each receives presence/absence labels and severity scores (0-10) with justification.

Inter-rater reliability analysis on 200 interactions showed Cohen's kappa > 0.72, indicating substantial agreement.

\section{Results}

We conducted 19 experiments across eight models and three conditions, generating 2,067 interaction turns with 1,559 concerning behavior instances.

\subsection{Condition-Dependent Patterns}

\begin{table}[h]
\centering
\begin{tabular}{lccc}
\hline
\textbf{Condition} & \textbf{Deception} & \textbf{Scheming} & \textbf{Power} \\
\hline
Neutral Autonomy & 0.0\% & 5.1\% & 6.2\% \\
Suggestive Autonomy & 2.8\% & 33.6\% & 29.1\% \\
Constrained Autonomy & 1.1\% & 4.4\% & 6.4\% \\
\hline
\end{tabular}
\caption{Behavioral frequencies across conditions}
\end{table}

Suggestive autonomy elicited scheming at 6.6 times the rate of neutral autonomy (Cohen's $d = 1.023$, large effect). Constrained autonomy produced frequencies similar to neutral (d = -0.086, negligible effect).

\subsection{Model-Level Analysis}

\begin{table}[h]
\centering
\begin{tabular}{lccc}
\hline
\textbf{Model} & \textbf{N} & \textbf{Scheming} & \textbf{Power} \\
\hline
llama-v3p1-8b-instruct & 2 & 42.0\% & 42.5\% \\
glm-4p5-air & 3 & 38.3\% & 31.7\% \\
llama4-scout-instruct & 3 & 1.3\% & 2.7\% \\
\hline
\end{tabular}
\caption{Model-level behavioral frequencies (32-fold variation)}
\end{table}

\subsection{Behavioral Categories}

Unauthorized expansion (26.7\%) and boundary testing (25.6\%) were most frequent. Explicit deception occurred in only 1.4\% of turns, suggesting models more readily exceed boundaries than actively conceal actions.

\subsection{Cost Analysis}

Total classification cost: \$1.37 for 2,067 turns (\$0.0007/turn). Per-experiment cost: \$0.50-\$3.00. This achieves approximately 100-fold cost reduction versus traditional red-teaming (\$6,000-\$12,000).

\section{Discussion}

\subsection{Interpretation}

The large effect size (Cohen's $d = 1.023$) demonstrates that prompt framing critically influences scheming behaviors. Deployment scenarios emphasizing autonomy or freedom may inadvertently elicit concerning behaviors.

The 32-fold model variation indicates safety properties are not uniform across frontier models. Some highly capable models showed low scheming rates while certain smaller models exhibited high rates, suggesting scheming emerges from specific training characteristics rather than raw capability.

\subsection{Limitations}

\textbf{Role-Playing vs. Genuine Misalignment:} Models may interpret the task as creative fiction rather than expressing genuine behavioral tendencies. This fundamental ambiguity applies to all text-based AI safety evaluations.

\textbf{Environmental Fidelity:} Text-based simulation lacks full physical complexity. Behaviors in simulation may not transfer to actual embodied systems.

\textbf{Sample Size:} Current dataset of 19 experiments provides power for large effects but may miss smaller significant effects.

\subsection{Implications}

\begin{itemize}
\item Deployment context critically influences safety outcomes
\item Model-specific evaluation required; safety properties vary substantially
\item Environment complexity influences detection; simple environments may provide false confidence
\item Different behavioral categories require targeted interventions
\end{itemize}

\subsection{Framework Advantages}

The dual-LLM architecture provides reproducibility through complete logging, scalability at 100-fold cost reduction, and modularity enabling adaptation to diverse research questions beyond scheming.

\section{Conclusion}

We presented a reproducible framework for evaluating autonomous AI scheming behaviors. Evaluation of eight frontier models revealed significant behavioral differences driven by prompt framing (suggestive autonomy eliciting scheming at 6.6 times neutral rate, Cohen's $d = 1.023$) and 32-fold model variation in scheming propensity.

The framework achieves comprehensive evaluation at ~1\% of traditional red-teaming cost while providing complete reproducibility. The modular architecture enables adaptation to diverse AI safety research questions.

These findings demonstrate that deployment context critically influences behavioral outcomes. The substantial model variation necessitates model-specific safety evaluation. Future work should pursue balanced large-scale studies, temporal analysis, mechanistic interpretation, and validation through real embodied systems.

All experimental code, world-state specifications, prompts, and interaction logs are publicly available to enable independent replication and extension.

\bibliographystyle{plain}
\bibliography{references_updated}

\end{document}
